\chapter{Трёхступенная Real-связывающая суперсвязность и граница полного квадрата}
\label{ch:v7-real-linking-superconnection-assembly-gate}

\section{Можно ли получить относительную кривизну квадратом}

Предыдущий гейт получил каноническую бимодульную кривизну
\begin{equation}
 \mathcal R_U(A)=C_t(A)U-UC_s(A)
 =AA^*U-UA^*A,
 \label{eq:v7-linking-assembly-relative-curvature}
\end{equation}
где сокращение фоновых членов следует из полярного переплетения. Оставался
вопрос, является ли эта формула настоящим блоком квадрата одного нечётного
оператора или только составным выражением, вставленным после вычисления.

Пусть
\begin{equation}
 H_0=\mathbb C^{11},\qquad
 H_1=\mathbb C^{11}\oplus\mathbb C^{10},\qquad
 H_2=\mathbb C^{10}.
 \label{eq:v7-linking-assembly-chain-spaces}
\end{equation}
Определим две линейные по $A$ стрелки
\begin{equation}
 B_0(A)=
 \begin{pmatrix}A^*U\\A\end{pmatrix}:H_0\to H_1,
 \qquad
 B_1(A)=
 \begin{pmatrix}A&-UA^*\end{pmatrix}:H_1\to H_2.
 \label{eq:v7-linking-assembly-arrows}
\end{equation}
Тогда прямое блочное умножение даёт
\begin{equation}
 \boxed{B_1(A)B_0(A)=AA^*U-UA^*A=\mathcal R_U(A).}
 \label{eq:v7-linking-assembly-factorization}
\end{equation}
Таким образом, знак разности не добавляется в действие: он является знаком
между двумя длины-два путями одного трёхступенного комплекса.

\section{Нечётный оператор и его квадрат}

На носителе
\begin{equation}
 \mathcal H_{\rm link}=H_0\oplus H_1\oplus H_2,
 \qquad \dim_{\mathbb C}\mathcal H_{\rm link}=42,
 \label{eq:v7-linking-assembly-carrier}
\end{equation}
введём
\begin{equation}
 d_A=
 \begin{pmatrix}
 0&0&0\\
 B_0(A)&0&0\\
 0&B_1(A)&0
 \end{pmatrix},
 \qquad
 \chi_{\rm link}=\operatorname{diag}(I_{11},-I_{21},I_{10}).
 \label{eq:v7-linking-assembly-differential}
\end{equation}
Имеем
\begin{equation}
 \{\chi_{\rm link},d_A\}=0,
 \qquad
 [\chi_{\rm link},d_A^2]=0,
 \label{eq:v7-linking-assembly-parity}
\end{equation}
и единственный ненулевой блок длины два равен
\begin{equation}
 (d_A^2)_{20}=\mathcal R_U(A).
 \label{eq:v7-linking-assembly-square-block}
\end{equation}
На фоне $A=A_0$ и в нуле этот блок исчезает с машинными остатками меньше
$5\cdot10^{-15}$. На ста случайных комплексных матрицах максимальный
остаток тождества~\eqref{eq:v7-linking-assembly-factorization} меньше
$3.2\cdot10^{-14}$.

Это положительный результат: относительная полярная кривизна действительно
является квадратом заранее типизированного нечётного оператора, а не
постфактум функционалом.

\section{Почему полный самосопряжённый квадрат не проходит}

Real-самосопряжённое завершение имеет вид
\begin{equation}
 Q_A=d_A+d_A^*.
 \label{eq:v7-linking-assembly-selfadjoint-operator}
\end{equation}
Однако $Q_A^2$ содержит не только крайний блок $B_1B_0$, но и диагональные
положительные блоки
\begin{equation}
 B_0^*B_0,\qquad
 B_0B_0^*+B_1^*B_1,\qquad
 B_1B_1^*.
 \label{eq:v7-linking-assembly-diagonal-blocks}
\end{equation}
Именно они восстанавливают уже запрещённые суммарные Gram-каналы.
Для полного действия
\begin{equation}
 \mathcal S_{\rm full}
 =\mathcal S_E+\frac12
 \|Q_A^2-Q_{A_0}^2\|_{\rm HS}^2
 \label{eq:v7-linking-assembly-full-action}
\end{equation}
нулевой гессиан равен
\begin{equation}
 (n_-,n_0,n_+)=(27,0,0),
 \qquad
 \lambda_{\rm heavy}^{\min}=-\frac{448}{5}.
 \label{eq:v7-linking-assembly-full-origin-fail}
\end{equation}
То есть полный квадрат делает неустойчивыми уже все исследуемые направления.

Даже обычные множители, связанные с Real-полуследом и удалением
ориентационного двойного счёта, не спасают конструкцию. Если $\alpha$
умножает полный связывающий гессиан, допустимое окно равно
\begin{equation}
 0\le\alpha<\frac1{15}.
 \label{eq:v7-linking-assembly-scale-window}
\end{equation}
Но последовательные стандартные кандидаты
\begin{equation}
 \alpha=1,\frac12,\frac14,\frac18
 \label{eq:v7-linking-assembly-standard-scales}
\end{equation}
дают соответственно сигнатуры
\begin{equation}
 (27,0,0),\ (27,0,0),\ (27,0,0),\ (21,0,6).
 \label{eq:v7-linking-assembly-standard-fails}
\end{equation}
Первый проходящий контроль $\alpha=1/16$ уже является новым секторным
ослаблением, а не следствием текущего Real-завершения.

Целевой вакуум полного квадрата остаётся строго устойчивым:
\begin{equation}
 (n_-,n_0,n_+)_{A_0}=(0,0,27),
 \qquad \lambda_{\min}=23.2624779769\ldots.
 \label{eq:v7-linking-assembly-full-vacuum}
\end{equation}
Провал относится не к насыщению, а к избыточному запуску из нуля.

\section{Сохранившийся блок степени два}

Если действие читает только канонический крайний блок
\begin{equation}
 \mathcal S_{(2)}
 =\mathcal S_E+\frac12\|(d_A^2)_{20}\|_{\rm HS}^2,
 \label{eq:v7-linking-assembly-degree-two-action}
\end{equation}
то оно совпадает с относительным полярным кандидатом предыдущего гейта и
даёт
\begin{align}
 (n_-,n_0,n_+)_{0}&=(7,0,20),
 &\lambda_{\rm heavy}^{\min}&=\frac{18}{5},
 \label{eq:v7-linking-assembly-degree-two-origin}\\
 (n_-,n_0,n_+)_{A_0}&=(0,0,27),
 &\lambda_{\min}&=3.9368554658\ldots.
 \label{eq:v7-linking-assembly-degree-two-vacuum}
\end{align}

Но операторная факторизация ещё не выводит право отбросить диагональные
блоки~\eqref{eq:v7-linking-assembly-diagonal-blocks}. Такое право должно
следовать из степени формы, представленного junk-идеала, отображающего
конуса или относительной когомологии комплекса. Без этого
\eqref{eq:v7-linking-assembly-degree-two-action} остаётся канонически
локализованным, но ещё не полным следом кривизны.

\begin{theorem}[Положительная факторизация и no-go полного квадрата]
Относительная полярная кривизна точно факторизуется как блок $d_A^2$
трёхступенного нечётного комплекса размерностей $11\to21\to10$. Однако
полная норма самосопряжённой кривизны $(d_A+d_A^*)^2$ неизбежно добавляет
диагональные Gram-блоки и даёт нулевую сигнатуру $(27,0,0)$. Правильный
селектор сохраняется только на компоненте длины два
$\operatorname{Hom}(H_0,H_2)$.
\end{theorem}

\begin{nogocriterion}[Запрет скрытой проекции полного квадрата]
Нельзя объявлять $\|(d_A^2)_{20}\|^2$ полным Real-кривизностным действием,
просто удалив диагональные блоки $Q_A^2$. Требуется вывести их исчезновение
или ортогональное отделение в представленном дифференциальном исчислении.
Нельзя также использовать проходящий множитель $1/16$ без независимого
следового происхождения.
\end{nogocriterion}

\begin{protocol}
\begin{enumerate}
 \item трёхступенный носитель: \textbf{$11\to21\to10$};
 \item общая комплексная размерность: \textbf{$42$};
 \item $d_A$ нечётен: \textbf{точно};
 \item $d_A^2$ чётен: \textbf{точно};
 \item $(d_A^2)_{20}=\mathcal R_U$: \textbf{точно};
 \item полный самосопряжённый квадрат: \textbf{провал $(27,0,0)$};
 \item критический полный вес: \textbf{$1/15$};
 \item стандартные Real/ориентационные множители: \textbf{не проходят};
 \item блок степени два: \textbf{проходит $(7,0,20)$};
 \item вакуум блока степени два: \textbf{$(0,0,27)$};
 \item представленная проекция степени два: \textbf{не выведена};
 \item следующий гейт: \textbf{quotient кривизны длины два}.
\end{enumerate}
\end{protocol}

Машинный сертификат сохранён в
\path{s2t_v7_real_linking_superconnection_assembly_gate_results.json}.

\begin{thebibliography}{9}
\bibitem{v7-linking-assembly-quillen}
D.~Quillen,
\emph{Superconnections and the Chern Character},
Topology 24 (1985), 89--95.
\bibitem{v7-linking-assembly-ackermann-tolksdorf}
T.~Ackermann, J.~Tolksdorf,
\emph{The Generalized Lichnerowicz Formula and Analysis of Dirac Operators},
J. Reine Angew. Math. 471 (1996), 23--42; arXiv:hep-th/9503153.
\bibitem{v7-linking-assembly-roepstorff}
G.~Roepstorff,
\emph{Superconnections and the Higgs Field},
arXiv:hep-th/9801040.
\end{thebibliography}